\documentclass[../../main.tex]{subfiles}

\begin{document}

\section{商拓扑}\label{section:ch02-商拓扑}

商拓扑与本章中研究过的几种拓扑不同，它并不是分析中讨论过的某些东西的自然推广。尽管如此，得到它也并不困难，一个办法就是从几何上引入，我们常常采用“切割与黏合”的手段来做出诸如曲面这样的几何对象。例如，环面（轮胎的表面）可以通过适当地“黏合”矩形的对边而得到，参见\cref{figure:ch02-商拓扑-22.1}。球面（球体的表面）可以通过把圆盘的整个边界捏成一点做出来，参见\cref{figure:ch02-商拓扑-22.2}。将这些做法正式地写出来，就包含了商拓扑的概念。

\par
\begin{figure}[htbp]
\centering
\includegraphics[width=0.6\textwidth]{ch02-fig22-1}
\caption{矩形黏合对边得到环面}\label{figure:ch02-商拓扑-22.1}
\end{figure}

\begin{definition}[商映射]\label{definition:ch02-商映射}
设 $X$ 和 $Y$ 都是拓扑空间，$p: X \rightarrow Y$ 是一个满射。如果 $Y$ 的子集 $U$ 是 $Y$ 的开集当且仅当 $p^{-1}(U)$ 是 $X$ 的一个开集，则称 $p$ 是一个商映射（quotient map）。
\end{definition}

\par
\begin{figure}[htbp]
\centering
\includegraphics[width=0.4\textwidth]{ch02-fig22-2}
\caption{圆盘边界捏成一点得到球面}\label{figure:ch02-商拓扑-22.2}
\end{figure}

这个条件比连续性强。有的数学家称之为“强连续性”。与它等价的条件是，$Y$ 的子集 $A$ 是 $Y$ 的闭集当且仅当 $p^{-1}(A)$ 是 $X$ 中的闭集。两者的等价性可由下面的式子得到：
\begin{align*}
f^{-1}(Y - B) = X - f^{-1}(B).
\end{align*}

描述商映射的另一个方法是：$X$ 的一个子集 $C$ 称为（关于满映射 $p: X \rightarrow Y$）是饱和的（saturated），如果 $C$ 包含每一个与之相交的 $p^{-1}(\{y\})$。或者这样说：$C$ 称为饱和集，如果 $C = p^{-1}(p(C))$。我们说 $p$ 是一个商映射，等于说 $p$ 是连续的，并且 $p$ 将 $X$ 的饱和的开集映成 $Y$ 的开集（$p$ 将 $X$ 的饱和的闭集映成 $Y$ 的闭集）。

两种特殊的商映射是开映射和闭映射。以前说过，一个映射 $f: X \rightarrow Y$ 称为开映射（open map），如果对于 $X$ 中每一个开集 $U$，$f(U)$ 是 $Y$ 中的开集。所谓 $f: X \rightarrow Y$ 是闭映射（closed map），是说对于 $X$ 的每一个闭集 $A$，$f(A)$ 是 $Y$ 中的闭集。根据定义直接可见，若 $p: X \rightarrow Y$ 是满的连续开映射或满的连续闭映射，则 $p$ 是一个商映射。此外，存在着既不是开映射又不是闭映射的商映射。（见习题 3。）

\begin{example}\label{example:ch02-商映射例1}
设 $X$ 是 $\mathbb{R}$ 的子空间 $[0, 1] \cup [2, 3]$，$Y$ 是 $\mathbb{R}$ 的子空间 $[0, 2]$。由下式
\begin{align*}
p(x) =
\begin{cases}
x, & \text{对于 } x \in [0, 1], \\
x - 1, & \text{对于 } x \in [2, 3]
\end{cases}
\end{align*}
定义的映射 $p: X \rightarrow Y$ 是连续的满的闭映射，但它不是开映射。因为 $X$ 中开集 $[0, 1]$ 的像不是 $Y$ 中的开集。

注意，若取 $X$ 的子集 $A$ 为 $[0, 1) \cup [2, 3]$，则由 $p$ 的限制所得的映射 $q: A \rightarrow Y$ 是一个连续的满射，但它不是商映射。这是由于 $[2, 3]$ 是一个开集并且相对于 $q$ 而言是一个饱和集，但它的像不是 $Y$ 的开集。
\end{example}

\begin{example}\label{example:ch02-商映射例2}
设 $\pi_{1}: \mathbb{R} \times \mathbb{R} \rightarrow \mathbb{R}$ 为到第一个坐标空间的投射，$\pi_{1}$ 是一个连续的满射。此外 $\pi_{1}$ 也是一个开映射。事实上，若 $U \times V$ 是 $\mathbb{R} \times \mathbb{R}$ 的一个基元素，则 $\pi_{1}(U \times V) = U$ 是 $\mathbb{R}$ 的一个开集。于是得到，$\pi_{1}$ 将 $\mathbb{R} \times \mathbb{R}$ 的开集映为 $\mathbb{R}$ 的开集。但是 $\pi_{1}$ 不是一个闭映射。例如，$\mathbb{R}$ 的子集
\begin{align*}
C = \{x \times y \mid xy = 1\}
\end{align*}
是一个闭集，但是 $\pi_{1}(C) = \mathbb{R} - \{0\}$ 不是 $\mathbb{R}$ 的闭集。
\end{example}

下面我们来说明如何用一个商映射来构造一个集合上的拓扑。

\begin{definition}[商拓扑]\label{definition:ch02-商拓扑}
设 $X$ 是一个空间，$A$ 是一个集合。$p: X \rightarrow A$ 是一个满射，则 $A$ 上恰好存在一个拓扑 $\mathcal{T}$，使 $p$ 为商映射。$\mathcal{T}$ 称为由 $p$ 导出的商拓扑（quotient topology）。

显然，这个拓扑 $\mathcal{T}$ 由 $A$ 中那些子集 $U$ 组成，这些 $U$ 满足条件：$p^{-1}(U)$ 在 $X$ 中是开的。
\end{definition}

容易验证，$\mathcal{T}$ 是一个拓扑。因为 $p^{-1}(\varnothing) = \varnothing$，$p^{-1}(A) = X$，所以 $\varnothing$ 与 $A$ 是开集。拓扑的另外两个条件可由下式推出，
\begin{align*}
p^{-1}\left(\bigcup_{\alpha \in J} U_{\alpha}\right) = \bigcup_{\alpha \in J} p^{-1}(U_{\alpha}), \\
p^{-1}\left(\bigcap_{i=1}^{n} U_{i}\right) = \bigcap_{i=1}^{n} p^{-1}(U_{i}).
\end{align*}

\par
\begin{figure}[htbp]
\centering
\includegraphics[width=0.4\textwidth]{ch02-fig22-3}
\caption{三点集的商拓扑}\label{figure:ch02-商拓扑-22.3}
\end{figure}

\begin{example}\label{example:ch02-商拓扑例3}
可以验证，用下式
\begin{align*}
p(x) =
\begin{cases}
a & \text{如果 } x > 0, \\
b & \text{如果 } x < 0, \\
c & \text{如果 } x = 0
\end{cases}
\end{align*}
定义的从实直线 $\mathbb{R}$ 到三点集 $A = \{a, b, c\}$ 的映射所导出的 $A$ 的商拓扑如\cref{figure:ch02-商拓扑-22.3}所示。
\end{example}

有一种经常要用到商拓扑的特殊情形如下：

\begin{definition}[商空间]\label{definition:ch02-商空间}
设 $X$ 是一个拓扑空间，$X^{*}$ 是 $X$ 的一个分拆，其成员是 $X$ 的两两无交的子集并且其成员的并为 $X$。设 $p: X \rightarrow X^{*}$ 是一个满射，其定义是将 $X$ 的每一点映到 $X^{*}$ 中包含这个点的唯一的一个元素。在用 $p$ 诱导出的商拓扑下，空间 $X^{*}$ 称为 $X$ 的一个商空间（quotient space）。
\end{definition}

当给定了 $X^{*}$，就有了一个 $X$ 上的等价关系，其等价类为 $X^{*}$ 的元素。人们可以把 $X^{*}$ 设想成“将每一个等价类中的所有元素黏合成一点”而得到的。因此，商空间 $X^{*}$ 也常称为 $X$ 的“黏合空间”或者“分解空间”。

可以用另一种方式来描述 $X^{*}$ 的拓扑。$X^{*}$ 的一个子集 $U$ 是一个等价类的集合，$p^{-1}(U)$ 恰好是属于 $U$ 的等价类的并。因此 $X^{*}$ 中的典型开集是一个等价类的族，它的并是 $X$ 中的开集。

\par
\begin{figure}[htbp]
\centering
\includegraphics[width=0.5\textwidth]{ch02-fig22-4}
\caption{闭单位球黏合边界得到 2-维球面}\label{figure:ch02-商拓扑-22.4}
\end{figure}

\begin{example}\label{example:ch02-商空间例4}
设 $X$ 为 $\mathbb{R}^{2}$ 中的闭单位球：
\begin{align*}
\{x \times y \mid x^{2} + y^{2} \leqslant 1\},
\end{align*}
设 $X^{*}$ 为 $X$ 的一个分拆，它由所有满足 $x^{2} + y^{2} < 1$ 的单点集 $\{x \times y\}$，以及集合 $S^{1} = \{x \times y \mid x^{2} + y^{2} = 1\}$ 组成。$X$ 中典型饱和开集如\cref{figure:ch02-商拓扑-22.4}中的阴影部分所示。可以证明：$X^{*}$ 同胚于 $\mathbb{R}^{3}$ 的一个叫做 2-维单位球面（unit 2-sphere）的子空间，其定义为
\begin{align*}
S^{2} = \{(x, y, z) \mid x^{2} + y^{2} + z^{2} = 1\}.
\end{align*}
\end{example}

\begin{example}\label{example:ch02-商空间例5}
设 $X$ 为矩形 $[0, 1] \times [0, 1]$。定义 $X$ 的一个分拆 $X^{*}$ 为：由所有单点集 $\{x \times y\}$（其中 $0 < x < 1$，$0 < y < 1$）、以下类型的两点集
\begin{align*}
\{x \times 0, x \times 1\}, \quad 0 < x < 1, \\
\{0 \times y, 1 \times y\}, \quad 0 < y < 1
\end{align*}
和四点集
\begin{align*}
\{0 \times 0, 0 \times 1, 1 \times 0, 1 \times 1\}
\end{align*}
组成。$X$ 中形如 $p^{-1}(U)$ 的几个典型饱和开集如\cref{figure:ch02-商拓扑-22.5}中的阴影部分所示。其中每一个都是 $X$ 的开集，并且等于等价类的并。
\end{example}

\par
\begin{figure}[htbp]
\centering
\includegraphics[width=0.6\textwidth]{ch02-fig22-5}
\caption{矩形中的饱和开集}\label{figure:ch02-商拓扑-22.5}
\end{figure}

这些集合中的每一个在 $p$ 下的像都是 $X^{*}$ 的开集，如\cref{figure:ch02-商拓扑-22.6}所示。对于 $X^{*}$ 的这种描述法，恰好就是用图示说明黏合矩形的边组成环面的一种数学表达。

\par
\begin{figure}[htbp]
\centering
\includegraphics[width=0.4\textwidth]{ch02-fig22-6}
\caption{商空间中的开集}\label{figure:ch02-商拓扑-22.6}
\end{figure}

现在研究一下商拓扑和商空间这两个概念与以前引进的一些概念之间的关系。值得注意的是，这关系不像我们所希望的那样简单。

我们已经发觉子空间的性质不够理想：若 $p: X \rightarrow Y$ 是一个商映射，$A$ 是 $X$ 的一个子空间，则由限制 $p$ 的定义域及值域而得到的映射 $q: A \rightarrow p(A)$ 就不一定还是商映射。然而，我们有以下定理：

\begin{theorem}[商映射的限制]\label{theorem:ch02-商映射限制}
设 $p: X \rightarrow Y$ 是一个商映射，$A$ 是 $X$ 的一个子空间，相对于 $p$ 饱和，$q: A \rightarrow p(A)$ 是 $p$ 的限制。

(1) 若 $A$ 是 $X$ 的一个开集或闭集，则 $q$ 是一个商映射。

(2) 若 $p$ 是一个开映射或闭映射，则 $q$ 是一个商映射。
\end{theorem}

\begin{proof}
第一步，我们首先验证以下两个等式：
\begin{align*}
q^{-1}(V) = p^{-1}(V) \quad \text{如果 } V \subset p(A); \\
p(U \cap A) = p(U) \cap p(A) \quad \text{如果 } U \subset X.
\end{align*}
为验证第一个等式，注意由于 $V \subset p(A)$ 和 $A$ 都是饱和集，$p^{-1}(V)$ 包含于 $A$。因此 $p^{-1}(V)$ 和 $q^{-1}(V)$ 都等于 $A$ 中经 $p$ 映入 $V$ 的点的集合。为验证第二个等式，注意对于 $X$ 的任何子集 $U$ 和 $A$，有以下包含关系
\begin{align*}
p(U \cap A) \subset p(U) \cap p(A).
\end{align*}
为了证明相反的包含关系，设 $y = p(u) = p(a)$，其中 $u \in U$ 且 $a \in A$。因为 $A$ 是一个饱和集，$A$ 包含集合 $p^{-1}(p(a))$，因此 $A$ 包含点 $u$。所以对于 $u \in U \cap A$，有 $y = p(u)$。

第二步，现假设 $A$ 是一个开集或者 $p$ 为开映射。对于 $p(A)$ 的给定的子集 $V$，我们来证明：若 $q^{-1}(V)$ 在 $A$ 中是开的，则 $V$ 在 $p(A)$ 中是开的。

首先设 $A$ 是一个开集。由于 $q^{-1}(V)$ 在 $A$ 中是开的，并且 $A$ 在 $X$ 中是开的，所以 $q^{-1}(V)$ 在 $X$ 中是开的。由于 $q^{-1}(V) = p^{-1}(V)$，而后者在 $X$ 中是开的，根据 $p$ 是一个商映射可见 $V$ 在 $Y$ 中是开的。从而，$V$ 在 $p(A)$ 中是开的。

再设 $p$ 是一个开映射。由于 $q^{-1}(V) = p^{-1}(V)$ 和 $q^{-1}(V)$ 是 $A$ 的开集，所以对于 $X$ 的某一个开集 $U$ 有 $p^{-1}(V) = U \cap A$。由于 $p$ 是一个满射，$p(p^{-1}(V)) = V$。所以
\begin{align*}
V = p(p^{-1}(V)) = p(U \cap A) = p(U) \cap p(A).
\end{align*}
由于 $p$ 是一个开映射，所以 $p(U)$ 是 $Y$ 的一个开集。因此 $V$ 在 $p(A)$ 中是开的。

第三步，$A$ 或者 $p$ 是闭的这种情形的证明可以通过在第二步中将“开”通通替换成“闭”而得到。
\end{proof}

现在我们考虑先前引入的几个概念。映射复合的性质较为理想。根据等式
\begin{align*}
p^{-1}(q^{-1}(U)) = (q \circ p)^{-1}(U),
\end{align*}
容易验证两个商映射的复合还是一个商映射。

此外，映射乘积的相关性质未如人意。两个商映射的笛卡儿积未必是一个商映射。见下面的例 7。为了得到相应的结论，需要对空间或者映射附加某些条件。对于空间所附加的一种条件是“局部紧致性”的概念，以后我们将会研究它。对于映射所附加的一种条件是：两个映射 $p$ 和 $q$ 都是商映射。容易验证此时 $p \times q$ 也是一个开映射，从而它也是一个商映射。

最后，Hausdorff 条件也不够理想。即便空间 $X$ 是 Hausdorff 的，我们却没有任何理由说商空间 $X^{*}$ 必须是 Hausdorff 的。有一个简单的条件能使得 $X^{*}$ 满足 $T_{1}$ 公理，这就是要求分拆 $X^{*}$ 中每一元素为 $X$ 的一个闭子集。保证 $X^{*}$ 成为 Hausdorff 空间的条件难于找到。这是一个有关商空间的较为棘手的问题。我们在本书中以后将多次回到这个问题。

有关商空间的研究中，或许最为重要的问题当属构造定义在商空间上的连续函数。现在我们就着手处理这一问题。我们曾有过一个判别法则，以确定一个到积空间中的映射 $f: Z \rightarrow \prod X_{\alpha}$ 的连续性。商空间理论中与之相应的判别法则便是确定一个定义从商空间出发的映射 $f: X^{*} \rightarrow Z$ 的连续性。我们有以下定理：

\begin{theorem}[商空间上映射的连续性]\label{theorem:ch02-商空间连续映射}
设 $p: X \rightarrow Y$ 是一个商映射。$Z$ 是一个空间，$g: X \rightarrow Z$ 是一个映射，对于每一个 $y \in Y$，$g$ 在每一个集合 $p^{-1}(\{y\})$ 上取常值。则 $g$ 诱导一个映射 $f: Y \rightarrow Z$，满足条件：$f \circ p = g$。映射 $f$ 是连续的当且仅当 $g$ 是连续的。$f$ 是一个商映射当且仅当 $g$ 是一个商映射。
\end{theorem}

\begin{proof}
对于每一个 $y \in Y$，集合 $g(p^{-1}(\{y\}))$ 是 $Z$ 中的一个单点集（因为 $g$ 在 $p^{-1}(\{y\})$ 上是常值）。如果用 $f(y)$ 表示这个点，那么我们便定义了一个映射 $f: Y \rightarrow Z$，使得对于每一个 $x \in X$，$f(p(x)) = g(x)$。如果 $f$ 连续，那么 $g = f \circ p$ 是连续的。反之，假定 $g$ 是连续的。令 $V$ 是 $Z$ 的一个开集，则 $g^{-1}(V)$ 是 $X$ 中的一个开集。然而 $g^{-1}(V) = p^{-1}(f^{-1}(V))$，由于 $p$ 是一个商映射，所以 $f^{-1}(V)$ 是 $Y$ 的一个开集。因此 $f$ 是连续的。

若 $f$ 为商映射，则 $g$ 作为两个商映射的复合也是一个商映射。反之，设 $g$ 为商映射。由于 $g$ 是一个满射，所以 $f$ 也是满的。设 $V$ 是 $Z$ 的一个子集，证明：如果 $f^{-1}(V)$ 在 $Y$ 中是开的，则 $V$ 在 $Z$ 中也是开的。由于 $p$ 是连续的，$p^{-1}(f^{-1}(V))$ 是 $X$ 中的一个开集。又由于 $p^{-1}(f^{-1}(V)) = g^{-1}(V)$，所以 $g^{-1}(V)$ 是 $X$ 的一个开集。由于 $g$ 是一个商映射，所以 $V$ 在 $Z$ 中是开的。
\end{proof}

\begin{corollary}\label{corollary:ch02-商空间同胚}
设 $g: X \rightarrow Z$ 是一个连续的满射。$X^{*}$ 为由下式
\begin{align*}
X^{*} = \{g^{-1}(\{z\}) \mid z \in Z\}
\end{align*}
定义的 $X$ 的子集族。在 $X^{*}$ 上取商拓扑。

(a) 则由映射 $g$ 诱导的映射 $f: X^{*} \rightarrow Z$ 是一一的连续映射，并且 $f$ 是同胚当且仅当 $g$ 是商映射。

(b) 若 $Z$ 是一个 Hausdorff 空间，则 $X^{*}$ 也是一个 Hausdorff 空间。
\end{corollary}

\begin{proof}
由前一个定理，$g$ 诱导了一个连续映射 $f: X^{*} \rightarrow Z$。显然 $f$ 是一一的。假设 $f$ 是一个同胚，则 $f$ 和投射 $p: X \rightarrow X^{*}$ 都是商映射，从而它们的复合 $g$ 也是商映射。反之，设 $g$ 是一个商映射。则根据前一个定理可见 $f$ 是商映射。由于 $f$ 是一一的，所以 $f$ 是一个同胚。

设 $Z$ 为 Hausdorff 空间。$X^{*}$ 中不同的点在 $f$ 下的像是不同的。因此存在无交的邻域 $U$ 和 $V$，于是 $f^{-1}(U)$ 与 $f^{-1}(V)$ 就是 $X^{*}$ 中给定两点的无交的邻域。
\end{proof}

\par
\begin{figure}[htbp]
\centering
\includegraphics[width=0.5\textwidth]{ch02-fig22-7}
\caption{将可数多个线段黏合于一端}\label{figure:ch02-商拓扑-22.7}
\end{figure}

\begin{example}\label{example:ch02-商空间例6}
设 $X$ 是由线段 $[0, 1] \times \{n\}$（其中 $n \in \mathbb{Z}_{+}$）之并构成的 $\mathbb{R}^{2}$ 的那个子空间，$Z$ 是由所有形如 $x \times (x/n)$ 的点构成的 $\mathbb{R}^{2}$ 的子空间，其中，$x \in [0, 1]$，$n \in \mathbb{Z}_{+}$。这时，$X$ 是可数多个无交线段之并，$Z$ 是有一公共端点的可数多个线段之并。参见\cref{figure:ch02-商拓扑-22.7}。

用公式 $g(x \times n) = x \times (x/n)$ 定义一个映射 $g: X \rightarrow Z$。$g$ 是一个满射，并且是连续的。以 $g^{-1}(\{z\})$ 为元素的商空间 $X^{*}$ 便是将 $X$ 的子集 $\{0\} \times \mathbb{Z}_{+}$ 黏合成一点所得到的空间。映射 $g$ 诱导出一个一一连续映射 $f: X^{*} \rightarrow Z$。但 $f$ 不是同胚。

为验证这一事实，我们只要证明 $g$ 不是商映射。考虑 $X$ 中点的序列 $x_{n} = (1/n) \times n$。由于集合 $A = \{x_{n}\}$ 没有极限点，因此它是 $X$ 中的一个闭集。它也是相对于 $g$ 而言的一个饱和集。另一方面，由于 $g(A)$ 是由所有形如 $z_{n} = (1/n) \times (1/n^{2})$ 的点所构成，所以 $g(A)$ 不是 $Z$ 的闭集，原点就是它的一个极限点。
\end{example}

\begin{example}\label{example:ch02-商空间例7}
两个商映射的积未必还是商映射。

在本节的习题中我们将给出一个涉及非 Hausdorff 空间的例子。我们在这里给出的例子所涉及空间都比较好。

令 $X = \mathbb{R}$，$X^{*}$ 是将 $X$ 的子集 $\mathbb{Z}_{+}$ 黏合成一点 $b$ 所形成的商空间。又令 $p: X \rightarrow X^{*}$ 为商映射。以 $\mathbb{Q}$ 表示 $\mathbb{R}$ 的所有有理数构成的子空间，设 $i: \mathbb{Q} \rightarrow \mathbb{Q}$ 为恒等映射。我们证明
\begin{align*}
p \times i: X \times \mathbb{Q} \longrightarrow X^{*} \times \mathbb{Q}
\end{align*}
不是商映射。

对于每一个 $n$，令 $c_{n} = \sqrt{2}/n$，并且考虑 $\mathbb{R}^{2}$ 中所有经过点 $n \times c_{n}$、斜率分别为 $1$ 和 $-1$ 的直线。设 $U_{n}$ 为 $X \times \mathbb{Q}$ 中位于两条直线上方或者两条直线下方并且位于两条垂线 $x = n - 1/4$ 与 $x = n + 1/4$ 之间的所有点的集合。则 $U_{n}$ 在 $X \times \mathbb{Q}$ 中是开的。由于 $c_{n}$ 不是有理数，$U_{n}$ 包含着集合 $\{n\} \times \mathbb{Q}$。参见\cref{figure:ch02-商拓扑-22.8}。
\end{example}

\par
\begin{figure}[htbp]
\centering
\includegraphics[width=0.6\textwidth]{ch02-fig22-8}
\caption{$p \times i$ 不是商映射}\label{figure:ch02-商拓扑-22.8}
\end{figure}

令 $U$ 为所有 $U_{n}$ 之并，则 $U$ 在 $X \times \mathbb{Q}$ 中是开的。对于每一个 $q \in \mathbb{Q}$，由于 $U$ 包含集合 $\mathbb{Z}_{+} \times \{q\}$，所以它是关于 $p \times i$ 的一个饱和集。我们假定 $U' = (p \times i)(U)$ 在 $X^{*} \times \mathbb{Q}$ 中是开的，并且由此导出矛盾。

特别地，由于 $U$ 包含着集合 $\mathbb{Z}_{+} \times 0$，所以集合 $U'$ 包含点 $b \times 0$。因此 $U'$ 包含着一个开集 $W \times I_{\delta}$，其中 $W$ 是 $b$ 在 $X^{*}$ 中的一个邻域，$I_{\delta}$ 由使得 $|y| < \delta$ 成立的所有有理数 $y$ 组成。这时便有，
\begin{align*}
p^{-1}(W) \times I_{\delta} \subset U.
\end{align*}

选取充分大的 $n$ 使得 $c_{n} < \delta$。由于 $p^{-1}(W)$ 在 $X$ 中是开的并且包含着 $\mathbb{Z}_{+}$，我们可以选取 $\varepsilon < 1/4$ 使得区间 $(n - \varepsilon, n + \varepsilon)$ 包含于 $p^{-1}(W)$。那么 $U$ 便包含着 $X \times \mathbb{Q}$ 的子集 $V = (n - \varepsilon, n + \varepsilon) \times I_{\delta}$。然而图形清晰地显示出：$V$ 的许多点 $x \times y$ 不在 $U$ 中！（这些点之一便是 $x \times y$，其中 $x = n + \frac{1}{2}\varepsilon$，$y$ 为满足 $|y - c_{n}| < \frac{1}{2}\varepsilon$ 的某一个有理数。）

\subsection*{习题}

\begin{exercise}\label{exercise:ch02-商拓扑-1}
详细验证例 3。
\end{exercise}

\begin{exercise}\label{exercise:ch02-商拓扑-2}
(a) 设 $p: X \rightarrow Y$ 是一个连续映射。证明：若存在一个连续映射 $f: Y \rightarrow X$ 使得 $p \circ f$ 等于 $Y$ 上的恒等映射，则 $p$ 为商映射。

(b) 如果 $A \subset X$，$X$ 到 $A$ 上的一个收缩（retraction）是一个连续映射 $r: X \rightarrow A$，并满足条件：对于每一个 $a \in A$，$r(a) = a$。证明：收缩是一个商映射。
\end{exercise}

\begin{exercise}\label{exercise:ch02-商拓扑-3}
设 $\pi_{1}: \mathbb{R} \times \mathbb{R} \rightarrow \mathbb{R}$ 是到第一个坐标空间的投射。令 $A$ 是由所有满足条件或者 $x \geqslant 0$ 或者 $y = 0$（或者二者同时成立）的点 $x \times y$ 所构成的 $\mathbb{R} \times \mathbb{R}$ 的那个子空间。$q: A \rightarrow \mathbb{R}$ 是 $\pi_{1}$ 的限制。证明：$q$ 是非开非闭的商映射。
\end{exercise}

\begin{exercise}\label{exercise:ch02-商拓扑-4}
(a) 在平面 $X = \mathbb{R}^{2}$ 中定义等价关系如下：
\begin{align*}
x_{0} \times y_{0} \sim x_{1} \times y_{1} \quad \text{如果 } x_{0} + y_{0}^{2} = x_{1} + y_{1}^{2}.
\end{align*}
令 $X^{*}$ 为相应的商空间，则 $X^{*}$ 同胚于一个你所熟悉的空间，它是什么空间？[提示：令 $g(x \times y) = x + y^{2}$。]

(b) 对于等价关系
\begin{align*}
x_{0} \times y_{0} \sim x_{1} \times y_{1} \quad \text{如果 } x_{0}^{2} + y_{0}^{2} = x_{1}^{2} + y_{1}^{2},
\end{align*}
讨论 (a) 小题中的问题。
\end{exercise}

\begin{exercise}\label{exercise:ch02-商拓扑-5}
设 $p: X \rightarrow Y$ 是一个开映射。证明：若 $A$ 为 $X$ 的一个开集，则 $p$ 的限制 $q: A \rightarrow p(A)$ 为开映射。
\end{exercise}

\begin{exercise}\label{exercise:ch02-商拓扑-6}
设 $\mathbb{R}_{K}$ 表示具有 $K$-拓扑的实直线（见第\cref{section:ch02-拓扑的基}节）。令 $Y$ 是将集合 $K$ 黏合成一点所得到的 $\mathbb{R}_{K}$ 的商空间，$p: \mathbb{R}_{K} \rightarrow Y$ 为商映射。

(a) 证明 $Y$ 满足 $T_{1}$ 公理，但不是一个 Hausdorff 空间。

(b) 证明 $p \times p: \mathbb{R}_{K} \times \mathbb{R}_{K} \rightarrow Y \times Y$ 不是商映射。[提示：对角线不是 $Y \times Y$ 的闭集，但它的原像在 $\mathbb{R}_{K} \times \mathbb{R}_{K}$ 中是闭的。]
\end{exercise}

\end{document}